\section{Les plateformes et services digitaux}
\label{les-plateformes-et-services-digitaux}

\subsection{La vue d\textquotesingle ensemble}
\label{la-vue-densemble-6}

\begin{enumerate}
 \item
  L\textquotesingle exécutif national peut recourir à l\textquotesingle utilisation de plateformes digitales qui soutiennent les opérations de la Fédération.
 \item
  Toutes les plateformes et tous les services digitaux qui nécessitent un abonnement mensuel ou annuel doivent être approuvés par le Conseil d\textquotesingle administration.

  \begin{enumerate}
   \item
    Les responsables d\textquotesingle activité sont exemptés de cette exigence pour les plateformes acquises pour leur.s activité.s respective.s.
  \end{enumerate}
\end{enumerate}

\subsection{L\textquotesingle espace de travail Google}
\label{lespace-de-travail-google}

\begin{enumerate}
 \item
  La Fédération utilise l\textquotesingle espace de travail Google comme outil central de production et de collaboration, y compris pour les réunions, le courrier électronique et le partage de fichiers.
 \item
  Les postes suivants doivent obtenir les privilèges d\textquotesingle un « super-administrateur.trice » :

  \begin{enumerate}
   \item
    Le/la président.e
   \item
    Le/la VPC
   \item
    Le/la VPFA
   \item
    Le/la commissaire à l\textquotesingle informatique
  \end{enumerate}
\end{enumerate}

\subsection{Flickr}
\label{flickr}

\begin{enumerate}
 \item
  La Fédération utilise Flickr comme plateforme centrale d\textquotesingle hébergement d\textquotesingle images pour les événements et les activités de la FCÉG.
 \item
  Le/la VPC est responsable du maintien et de la gestion de cette plateforme.
\end{enumerate}

\subsection{Squarespace}
\label{squarespace}

\begin{enumerate}
 \item
  La Fédération utilise Squarespace comme plateforme d\textquotesingle hébergement de site web et de gestion de contenu.
 \item
  Le/la VPC est responsable du maintien et de la gestion de cette plateforme.
\end{enumerate}

\subsection{Les comptes de courriel des activités}
\label{les-comptes-de-courriel-des-activites}

\subsubsection{La vue d\textquotesingle ensemble}
\label{la-vue-densemble-courriel-activites}

\begin{enumerate}
 \item
  La FCÉG fournira des comptes Google Workspace aux membres du comité organisateur (CO) d\textquotesingle une activité.
 \item
  Chaque membre du CO se verra attribuer un compte de courriel individuel, afin que chaque membre n\textquotesingle ait besoin que d\textquotesingle un seul compte pour l\textquotesingle ensemble de ses activités liées à la Fédération.
 \item
  Le partage d\textquotesingle un même compte de connexion par plusieurs personnes est interdit, car cela représente un risque pour la sécurité des technologies de l\textquotesingle information et rend l\textquotesingle attribution des modifications de documents peu claire.
\end{enumerate}

\subsubsection{Les comptes de courriel individuels}
\label{les-comptes-de-courriel-individuels}

\begin{enumerate}
 \item
  Les membres du CO peuvent se voir attribuer un compte de courriel personnel dans le format suivant :
  \begin{enumerate}
   \item
    \texttt{prenom@activite\{annee\}.fceg.ca} (p.\ ex.\ \texttt{alice@conf20XX.fceg.ca}); ou
   \item
    \texttt{prenom.initiale@activite\{annee\}.fceg.ca} ou \texttt{prenom.nom@activite\{annee\}.fceg.ca}, dans le cas où deux membres du CO ou plus partagent le même prénom.
  \end{enumerate}
\end{enumerate}

\subsubsection{Les comptes de courriel de rôle}
\label{les-comptes-de-courriel-de-role}

\begin{enumerate}
 \item
  Les comptes de courriel de rôle (p.\ ex.\ \texttt{finances@activite\{annee\}.fceg.ca}) existent uniquement en tant que boîtes de réception déléguées. La connexion directe à ces comptes ne sera pas fournie; l\textquotesingle accès est plutôt accordé en déléguant la boîte de réception aux comptes individuels des membres du CO concernés.

  \begin{enumerate}
   \item
    Cela permet à plusieurs personnes partageant un même rôle de gérer une boîte de réception commune tout en conservant des comptes individuels distincts.
   \item
    Cela permet également à une personne de gérer plusieurs boîtes de réception de rôle depuis son compte individuel.
  \end{enumerate}

 \item
  Les comptes de rôle ne peuvent pas être utilisés comme comptes Google Drive ni pour tout autre stockage Google Workspace au-delà du courriel.
\end{enumerate}

\subsubsection{Les options de domaine de courriel des activités}
\label{les-options-de-domaine-de-courriel-des-activites}

\begin{enumerate}
 \item
  Les activités peuvent configurer leurs comptes de courriel selon l\textquotesingle une des structures de domaine suivantes :

  \begin{enumerate}
   \item
    \textbf{Comptes liés à l\textquotesingle année} (p.\ ex.\ \texttt{@activite\{annee\}.fceg.ca}) : Les comptes sont propres à une année de conférence donnée. L\textquotesingle activité peut accéder à ces comptes en tout temps.
   \item
    \textbf{Comptes liés à l\textquotesingle activité} (p.\ ex.\ \texttt{@activite.fceg.ca}) : Les comptes sont associés au nom de l\textquotesingle activité d\textquotesingle une année à l\textquotesingle autre. L\textquotesingle accès à ces comptes est soumis aux conditions suivantes :
    \begin{enumerate}
     \item
      L\textquotesingle activité doit attendre que la conférence de l\textquotesingle année précédente ait eu lieu avant d\textquotesingle obtenir l\textquotesingle accès, que la conférence de l\textquotesingle année précédente ait utilisé ou non des comptes liés à l\textquotesingle activité.
     \item
      Si la conférence de l\textquotesingle année précédente utilisait des comptes liés à l\textquotesingle activité, le CO entrant doit attendre le plus tôt des deux moments suivants : la préparation des comptes pour la transition, ou soixante (60) jours après la tenue de la conférence précédente.
    \end{enumerate}
  \end{enumerate}
\end{enumerate}

\subsection{L\textquotesingle organisation des documents des activités}
\label{lorganisation-des-documents-des-activites}

\begin{enumerate}
 \item
  Tous les fichiers créés dans le cadre de la planification d\textquotesingle une activité doivent être stockés dans un Drive partagé désigné pour cette activité au sein de l\textquotesingle espace de travail Google de la Fédération.
 \item
  Le/la vice-président.e aux services doit se voir accorder l\textquotesingle accès en tant que gestionnaire du Drive partagé de l\textquotesingle activité.
 \item
  Les documents ne doivent pas être stockés dans des comptes Google Drive personnels, des services de stockage en nuage externes ou tout autre stockage en dehors de l\textquotesingle espace de travail Google de la Fédération.
 \item
  Note : Les Drives partagés de Google Workspace ne prennent pas en charge les formulaires nécessitant le téléversement de fichiers. Un raccourci vers un dossier contenant ces formulaires doit être inclus dans le Drive partagé plutôt que d\textquotesingle y stocker les formulaires directement.
\end{enumerate}

\subsection{Discord}
\label{discord}

\subsubsection{Le serveur Discord public}
\label{le-serveur-discord-public}

\begin{enumerate}
 \item
  La Fédération maintient un serveur Discord public ouvert aux étudiant.e.s, aux président.e.s d\textquotesingle associations étudiantes en génie et aux VP Externes à travers le pays.
 \item
  Le/La vice-président.e aux communications est propriétaire du serveur Discord public à l\textquotesingle aide d\textquotesingle un compte lié à l\textquotesingle adresse courriel \texttt{vpc@cfes.ca}.
 \item
  Le/La commissaire à l\textquotesingle informatique peut se voir accorder l\textquotesingle accès administrateur pour soutenir le maintien et la gestion technologique du serveur.
 \item
  Le serveur Discord public peut être utilisé pour :

  \begin{enumerate}
   \item
    Faciliter les élections partielles;
   \item
    Solliciter la rétroaction des leaders étudiants;
   \item
    Offrir un espace de communication aux leaders étudiants à travers le pays; et
   \item
    Accueillir les heures de bureau de l\textquotesingle exécutif national.
  \end{enumerate}

\end{enumerate}

\subsubsection{Le serveur Discord interne}
\label{le-serveur-discord-interne}

\begin{enumerate}
 \item
  La Fédération maintient un serveur Discord interne destiné à l\textquotesingle équipe de la FCÉG.
 \item
  Le/La président.e est propriétaire du serveur Discord interne à l\textquotesingle aide d\textquotesingle un compte lié à l\textquotesingle adresse courriel \texttt{president@cfes.ca}.
 \item
  Le serveur Discord interne peut être utilisé pour les discussions au sein de l\textquotesingle exécutif national, du Conseil d\textquotesingle administration et de l\textquotesingle équipe de la FCÉG dans son ensemble concernant les affaires de la Fédération.
\end{enumerate}

\subsection{Les domaines}
\label{les-domaines}

\begin{enumerate}
 \item
  La Fédération utilise les domaines suivants pour ses sites web :

  \begin{enumerate}
   \item
    cfes-fceg.ca

    \begin{enumerate}
     \item
      Il s\textquotesingle agit de l\textquotesingle adresse principale du site web de la Fédération.
     \item
      Ce domaine sera désigné comme domaine alias de l\textquotesingle utilisateur pour cfes.ca dans l\textquotesingle espace de travail Google.
    \end{enumerate}
   \item
    cfes.ca

    \begin{enumerate}
     \item
      Il servira de domaine principal pour l\textquotesingle espace de travail Google et tous les comptes de courriel.
     \item
      L\textquotesingle adresse de ce site web sera redirigée vers le site cfes-fceg.ca.
    \end{enumerate}
   \item
    fceg.ca

    \begin{enumerate}
     \item
      L\textquotesingle adresse de ce site web sera redirigée vers le site cfes-fceg.ca.
     \item
      Ce domaine sera désigné comme domaine d\textquotesingle alias d\textquotesingle utilisateur pour le site cfes.ca dans l\textquotesingle espace de travail de Google.
    \end{enumerate}
  \end{enumerate}
\end{enumerate}

